\section{Sampling can inflate the dynamic PR integral}
\label{sec:sampling-inflation}
Fix the original dynamic VUS-PR contract for $n$ binary labels with both
classes present, $0<P_0<n$, finite scores, and a nonnegative integer
maximum width $L$. The implementation evaluates the scores after its
documented binary64 conversion. At width $w$, original positives
have unit weight and each interval adds a square-root ramp
$\sqrt{1-d/w}$ at distances $1\leq d\leq\lfloor w/2\rfloor$;
overlapping contributions are summed and clipped at one. At a complete
score threshold, let $C$ be the predicted count, $A$ the admitted original
positive count, $B$ the admitted buffer mass, $P_0$ the total original
positive count, and $h/K_w$ the fraction of expanded components hit.
Overlapping expanded intervals merge, whereas adjacent nonoverlapping
intervals remain distinct. The coordinates are
\[
 T=A+B,\qquad P=P_0+B/2,\qquad
 R=\min(T/P,1)h/K_w,\qquad \pi=T/C.
\]
The exhaustive PR area is $\sum_j(R_j-R_{j-1})\pi_j$, beginning at
$R_0=0$ and admitting each distinct global score tie completely. VUS-PR
averages this area over every integer width $0,\ldots,L$.

The sampled implementation retains descending zero-based ranks
\texttt{linspace(0,n-1,H).astype(int)}, with $H=250$ by default
in both the official VUS optimized API~\citep{officialvuscode} and
the pinned fast-VUS source~\citep{klein2026}. It uses the same PR
integral on those retained states. Sampling error is different from a
change of metric semantics or an implementation defect: the upstream
routine correctly computes its configured sampled target.

The important feature is that precision can increase \emph{after} recall
has already reached one. Additional buffer mass then contributes zero
exhaustive area, but a coarse threshold can incorrectly assign that later
precision to a recall jump which occurred earlier. Binary average
precision cannot add positive mass after every original positive has
already been admitted. The following construction isolates the additional
soft-label mechanism.

\begin{proposition}[A parameterized sampling family]
\label{prop:overshoot-family}
Let $H\geq2$, $1\leq m<q$, $n=1+(H-1)q$, and
$\lfloor L/2\rfloor<n-m$. Put a single original positive at temporal
position zero. Assign distinct descending scores first to the last $m$
positions, then to position zero, and then to positions
$1,\ldots,n-m-1$ in temporal order. The exhaustive VUS-PR equals
$E=1/(m+1)$, whereas its $H$-threshold sampled value is
\begin{equation}
 S=\frac{1+(L+1)^{-1}\sum_{w=0}^L B_w}{q+1},\qquad
 B_w=\sum_{d=1}^{\min(q-m,\lfloor w/2\rfloor)}\sqrt{1-d/w},
 \label{eq:overshoot-value}
\end{equation}
where the empty sum at $w=0$ is zero. The sign of $S-E$ depends on the
parameters; the overestimation regimes below are separate conclusions.
\end{proposition}
\begin{proof}
The first $m$ positions are outside every evaluated support. At actual
rank $m+1$, admitting the sole original positive makes both coverage and
clipped recall one. The entire exhaustive recall increase therefore has
precision $1/(m+1)$, independently of width. Later recall increments vanish.
The retained one-based ranks are exactly $1,1+q,\ldots,n$. Recall at the
first is zero; recall at the second is one, and the second includes the
original plus the nearest $q-m$ positions. Its precision is
$(1+B_w)/(q+1)$, proving the formula after averaging widths.
\end{proof}

\begin{corollary}[Large finite overestimation]
\label{cor:large-overshoot}
For $H=250$ there is a finite member of
Proposition~\ref{prop:overshoot-family} with $S-E>0.98$.
\end{corollary}
\begin{proof}
For $L\geq2q$, retain only widths $2q,\ldots,L$ and use
$\sqrt{1-x}\geq1-x$ to obtain
\[
 S\geq\frac{(L-2q+1)(q-m+1)
       -\frac{(q-m)(q-m+1)}2\sum_{w=2q}^Lw^{-1}}
       {(L+1)(q+1)}.
\]
Take $q=2{,}000{,}000$, $m=1000$, $n=498{,}000{,}001$ and
$L=498{,}000{,}000$. The harmonic tail is less than
$\log(L/(2q-1))<5$; the latter inequality follows, for example, from
the degree-seven positive Taylor lower bound for $\exp(5)$.
Thus $S>977516488501501/996000500000001$. Exact rational subtraction
of $E=1/1001$ exceeds $98/100$.
\end{proof}
This finite statement was verified symbolically, without allocating the
large input or executing an evaluator on it. Its maximum width is $n-1$;
it is not a claim about a typical small tolerance width. More generally,
with fixed $H>3$, $L=(H-1)q$, $m\to\infty$ and $m/q\to0$,
\[
 \liminf S\geq1-\frac2{H-1}
                   -\frac{\log((H-1)/2)}{2(H-1)},\qquad E\to0.
\]
This asymptotic family is a real-arithmetic statement, rather than an
unbounded supply of distinct binary64 scores.

\paragraph{Large selection regret at ordinary fixed widths.}
For a second score vector on the same labels, place a nearby buffer
position $d$ first and the original second. Fill the remaining $q-1$
positions through rank $q+1$ with far outside-support negatives; rank
everything else later. Sufficient feasibility conditions are
$n-q+1>\lfloor L/2\rfloor$ and $1\leq d<n-q+1$.
At $w<2d$ the exhaustive and sampled areas are respectively
$1/2$ and $1/(q+1)$. Otherwise, writing
$a=\sqrt{1-d/w}$ and $r=a/(1+a/2)$ gives
\[
 E_B(w)=ra+(1-r)(1+a)/2,\quad
 S_B(w)=ra+(1-r)(1+a)/(q+1).
\]
Table~\ref{tab:constructed-reversal} reports complete width averages.
All scores are distinct finite integers. Closed forms, independent dense
evaluation, compiled exhaustive evaluation, and unchanged pinned upstream
calls at $H=250$ and $H=n$ agree numerically. These are exploratory
constructed witnesses of possibility, not estimates of prevalence.

\begin{table}[t]
\centering\small
\begin{tabular}{rrrrrrrr}
\toprule
$L$ & $n$ & $(q,m,d)$ & $E_A$ & $S_A$ & $E_B$ & $S_B$ & regret\\
\midrule
30 & 2491 & $(10,3,5)$ & .250000 & .506412 & .757142 & .418181 & .507142\\
100 & 4981 & $(20,4,5)$ & .200000 & .624743 & .900580 & .569008 & .700580\\
300 & 9961 & $(40,6,1)$ & .142857 & .706289 & .989249 & .668767 & .846392\\
\bottomrule
\end{tabular}
\caption{Same-label reversals with $H=250$. Sampling selects A; exhaustive
evaluation selects B. Regret is $E_B-E_A$, measured in exhaustive VUS-PR
units and not in an operational detection utility.}
\label{tab:constructed-reversal}
\end{table}

\section{Sound bounds from retained threshold states}
\label{sec:threshold-certificate}
An approximation can be accompanied by a checkable interval instead of
being treated as an exact scalar. The general idea of bounding numerical
quadrature is established; our result supplies a bound for this dynamic,
clipped-recall, soft-label target without assuming monotone precision.

\begin{theorem}[Mass-envelope PR enclosure]
\label{thm:threshold-enclosure}
At one width, consider two consecutive retained complete-threshold states
$(c_a,t_a,r_a)$ and $(c_b,t_b,r_b)$, including an initial zero state and
discarding repeated states. Put $\delta=t_b-t_a$ and
\[
 \ell=\frac{t_a}{c_b-\delta},\qquad
 u=\frac{t_b}{c_a+\delta}.
\]
If $c_b-\delta=0$, set $\ell=1$; a zero-mass block contributes zero.
The exhaustive PR area inside this block is between
\begin{equation}
 \int_{r_a}^{r_b}\max(\ell,P_0r/c_b)\,dr
 \quad\text{and}\quad (r_b-r_a)u.
 \label{eq:block-bounds}
\end{equation}
Summing these bounds over blocks and averaging widths encloses exhaustive
VUS-PR. In particular, disjoint model intervals certify their exact order.
\end{theorem}
\begin{proof}
Every temporal weight lies in $[0,1]$. At an intermediate predicted count
$c$, its true mass $t$ satisfies
\[
 \max(t_a,t_b-c_b+c)\leq t\leq\min(t_b,t_a+c-c_a).
\]
Dividing by $c$ and optimizing the two continuous piecewise envelopes
gives $\ell\leq t/c\leq u$. Optimizing over continuous counts is safe
for integer counts and complete ties because it enlarges the feasible
set. If the lower denominator vanishes, all admitted weights are one.

Recall is nondecreasing along descending scores: $(A+B)/(P_0+B/2)$
increases in $A$ and $B$, with derivative in $B$ proportional to
$P_0-A/2>0$. Clipping and multiplication by nonnegative, nondecreasing
coverage preserve this property. Thus all recall increments are nonnegative.
Also $T\geq P_0R$, and $C\leq c_b$, so contributing precision is at
least $P_0R/c_b$. The exact right-step sum dominates the integral of the
nondecreasing lower envelope; its upper bound follows by summing recall
increments times $u$.
\end{proof}

The bound is elementary and can be loose. Refining a retained grid
monotonically tightens the real-arithmetic enclosure: $c-t$ is
nondecreasing, so each subblock has a lower mass envelope at least as
large, an upper envelope no larger, and a no-smaller linear lower envelope.
If consecutive global distinct-score ordinals establish that no valid
threshold was omitted, use the observed exact right-step term instead.
A complete tie can admit many positions without containing any valid
intermediate threshold.

\paragraph{An information limit.}
Endpoint-only bounds cannot always be narrow. Fix $H\geq3$ and $L$;
choose one original interval of $m$ points and a sampling gap $q\gg m+L$.
Two score orders can put those originals immediately after the first
off-support negative, or at the end of the same gap, while giving identical
$(C,A,B,h)$ at every retained threshold and every width. Their exhaustive
areas are obtained by placing the entire maximum-width halo in a common
tail after that gap and filling the remaining gap with off-support points.
Recall saturates before the common halo enters. The exhaustive
areas are respectively
\[
 \frac1m\sum_{j=1}^m\frac{j}{j+1},\qquad
 \frac1m\sum_{j=1}^m\frac{j}{q+1-m+j}.
\]
The first tends to one and the second to zero as $m\to\infty$ and
$q/m\to\infty$. A bound valid for every score ordering consistent with
that retained transcript must include both. This is an ordinary ranking
information-loss argument; it does not lower-bound methods that inspect
omitted support scores or adapt their thresholds.

\section{Integer-directed realization}
\label{sec:integer-certificate}
Theorem~\ref{thm:threshold-enclosure} concerns real arithmetic. Adding a
floating-point tolerance to a numerical check would not implement that
guarantee. We instead construct dyadic enclosures using integer square
roots and directed integer division.

Let $D=2^{32}$ be the weight scale. For a single ramp at distance $d$,
set $z=(w-d)D^2$ and $q=\operatorname{isqrt}(\lfloor z/w\rfloor)$.
Then $q/D\leq\sqrt{1-d/w}\leq(q+1)/D$, with exact upper numerator
$q$ when $q^2w=z$. Original positives and points covered by two ramps have
exact numerator $D$; inactive points have zero. Since every active ramp
is at least $1/\sqrt2$, two ramps indeed clip to one. For $n<2^{31}$,
all compiled signed-64-bit cumulative buffer sums are below $2^{63}$.
Nonlinear calculations and final sums use arbitrary-precision Python
integers; public NumPy integer parameters are normalized before them.

For an endpoint with exact original count $A$, buffer brackets
$[B_-,B_+]/D$, and exact hit count $h$, use
\[
 f(B)=\frac{h\min\{2(AD+B),\,2P_0D+B\}}
              {K_w(2P_0D+B)}.
\]
Monotonicity gives $f(B_-)\leq R\leq f(B_+)$. With output scale
$S=2^{56}$, floor the lower product $Sf(B_-)$ and ceil the upper
product $Sf(B_+)$. Both endpoint precision envelopes can be rewritten
as functions nondecreasing in each total-mass argument:
\[
 \ell(t_a,t_b)=\frac{t_a}{c_b-t_b+t_a},\qquad
 u(t_a,t_b)=\frac{t_b}{c_a-t_a+t_b}.
\]
Substitute both lower mass endpoints into $\ell$ and both upper
endpoints into $u$. Combine these rational precision bounds with directed
recall differences. When $R_a^+<R_b^-$, the linear lower refinement is
integrated analytically over $[R_a^+,R_b^-]$, a subset of the actual recall
interval; otherwise that refinement contributes zero. Floor each
lower block numerator, ceil each upper numerator, then apply the same
directions after width averaging. No irrational or floating-point
operation enters this chain of bounds.

The returned certificate retains integer endpoints and scale. Strict
ordering compares integer cross-products, while float displays are
rounded outward and are not used for the decision. The certified target
is ideal real-arithmetic VUS-PR, not a promise to contain the exact bit
pattern of every floating-point summation implementation. An application
requiring nested refinement intervals must intersect independently rounded
bounds. This implementation is not an adaptive refinement policy.

An independent standard-library calculation checked 14,400 enclosures
against exact rational evaluations at deliberately coarse weight/output
precisions of three/twelve bits, and 12,500 root brackets by squared-integer
inequalities. The production implementation is additionally checked against
an independent 90-digit direct-ramp oracle, complete ties, boundary and
overlapping intervals, saturated recall, and signed/unsigned NumPy integer
parameters. These checks exercise the implementation; the universal
guarantee follows from the inequalities and accumulator bounds above.
